\documentclass[a4paper,11pt,openany]{ltjsbook}
\usepackage{luatexja}
\usepackage{luatexja-fontspec}
\usepackage{amsmath,amssymb,amsfonts,amsthm}
\usepackage{geometry}
\geometry{top=24mm,bottom=24mm,left=24mm,right=24mm,footskip=10mm}
\usepackage{tikz}
\usepackage{tcolorbox}
\tcbuselibrary{skins,breakable}
\usepackage{hyperref}
\usepackage{fancyhdr}

\definecolor{primaryblue}{RGB}{20, 60, 110}
\definecolor{softblue}{RGB}{238, 244, 252}
\definecolor{darkslate}{RGB}{35, 45, 60}

\hypersetup{
    colorlinks=true,
    linkcolor=primaryblue,
    citecolor=darkslate,
    pdftitle={古典で学び直すラウス・ボール数学史入門},
    pdfauthor={教育情報分析研究会}
}


\newcommand{\opp}{\text{opp}}
\newcommand{\adj}{\text{adj}}
\newcommand{\hyp}{\text{hyp}}
\providecommand{\cosec}{\operatorname{cosec}}
\newcommand{\TEntry}[2][]{#2}
\newcommand{\Reason}[2][]{#2}
\newcommand{\BAR}{\rule{0.25in}{1pt}}
\newcommand{\gcm}{\operatorname{gcm}}

\pagestyle{fancy}
\fancyhf{}
\fancyhead[LE,RO]{\thepage}
\fancyhead[RE]{\leftmark}
\fancyhead[LO]{\rightmark}
\setlength{\headheight}{14pt}
\renewcommand{\headrulewidth}{0.4pt}

\title{
    {\huge\bfseries\gtfamily 古典で学び直すラウス・ボール数学史入門}\\[14pt]
    {\large 底本：\textit{A Short Account of the History of Mathematics} (Project Gutenberg eBook \#31246)}
}
\author{\textbf{教育情報分析研究会} 編著}
\date{2026年10月}

\begin{document}
\maketitle
\frontmatter
\tableofcontents
\clearpage
\mainmatter

\chapter{Section 1}

\#\#\# INTRODUCTION

この電子書籍は、Project Gutenbergによって公開されており、誰でも無料で利用できます。あなたはそれをコピーしたり、配布したり、再利用したりすることができます。これはProject Gutenbergライセンスに従っており、詳しくはwww.gutenberg.orgをご覧ください。

---

\#\#\#\# 短い数学の歴史概要

このページには番号が付けられていません。
ページ番号付けはアルファベット順に開始されます。

**数学の短い歴史概要**

このページには番号が付けられていません。

中心部分
\vspace*{\stretch{1}}
\begin{center}
\vspace*{\stretch{1}}
{\Large 短い歴史概要}

\vspace{10mm}

{\small の}

\vspace{10mm}

{\Huge 数学の}

\vspace{10mm}

{\small 歴史について}

\vspace{30mm}

{\small によって}

\vspace{10mm}

W. W. ローズ・ボール

{\scriptsize タンティリティコールègeのフェロー}

\vspace*{\stretch{2}}
ドーバー出版株式会社

ニューヨーク
\end{center}

新しいドーバー版は、1960年に最初に出版され、著者の最後の修正版である---1908年に出版された4番目の編集版---を未変更で再発行しています。

\vspace*{\stretch{2}}

**国際標準書籍番号：0-486-20630-0**

**図書館協会カタログカード番号：60-3187**

製造はアメリカ合衆国の製紙所で行われました。

ドーバー出版株式会社

180 Varick Street

ニューヨーク、N. Y. 10014

\clearpage
このページには番号が付けられていません。
\vspace*{\stretch{1}}
\begin{center}
\small
電子書籍を生成したGreg Lindahl, Viv, Juliet Sutherland, Nigel Blowerおよびhttp://www.pgdp.netのオンライン分散校正チーム
\end{center}

---

\#\#\#\# INTRODUCTION

この序文の末尾にある目次から、整理の方法が把握できるでしょう。概要は次の通りです。

最初の章（第1章）では、エジプト人およびフェニキア人の数学に関する既知の知識について簡潔に述べています。これはギリシャの影響下における数学の歴史への導入となります。その後の歴史は3つの時期に分かれています：

1. **まず、ギリシャの影響下にある時期（第2章から第7章）**：この時期には、多くの重要な数学的発見がなされました。
2. **次に、中世とルネサンス期（第8章から第13章）**：この時期は、数学がより広範囲にわたる影響を及ぼし始めました。
3. **最後に、近代期（第14章から第19章）**：この時期には、数学の発展が急速に進展しました。

これらの時期の数学について議論する際には、歴史における主要な出来事のみを述べるという制限のもとで、しばしば比較的小さな影響力を持つ人物や作品については無視してしまいました。おそらく、記述された数学者たちの発見が過大評価されている理由は、それ以前に先駆的な著作を書いた未引用の著者の存在によるものかもしれません。

また、実用天文学者や物理学者による結果については、それらが数学的に興味深い場合を除き、一般的には参照していません。引用した結果は通常、現代の記号を使用していますので、読者はその内容が他の著者に言及されているものと同じであることを覚えておくべきです。ただし証明は時としてより便利で親しみやすい言語に翻訳されている場合があります。

大半の記述は既存の歴史書や回想録からまとめたものです。議論される内容は、数学の発展を歴史的にまとめたもので、その発展に大きく貢献した人物たちの生涯と発見を例示して説明しています。

---

この章では、まずエジプト人およびフェニキア人の数学について触れることから始めましょう。彼らは古代世界において重要な役割を果たし、その知識がギリシャの数学に大きな影響を与えました。エジプト人は分数や幾何学的な問題を解く方法を発展させ、フェニキア人は計算と記録のためのアルファベット文字を使用しました。

\#\#\# エジプト人の数学

エジプト人たちは約4000年前から数学を用いて農業や建築などに活用していました。彼らは分数を扱う方法を発展させ、特に2分の1, 3分の1, 7分の1などの単純な分数については具体的な表記法を持っています。また、彼らは長方形の面積や三角形の面積を計算するための公式も開発しました。

\#\#\# フェニキア人の数学

フェニキア人は約3000年前からアルファベット文字を使用し始めました。これは数値を表すためにも使用され、これが後のローマ数字やヒンディー・アラビア数字の発展に影響を与えました。

これらの知識はギリシャの数学家たちによって継承され、発展させられることになります。次回は、その歴史的な背景と重要な人物について詳しく説明します。

---

この解説講義では、エジプト人およびフェニキア人の数学を紹介しました。彼らの知識がギリシャの数学にどのように影響を与えたかを理解することで、数学の発展を歴史的に俯瞰することができます。次回は、その続きとしてギリシャの数学家たちについて詳しく説明します。

---

以上のように、エジプト人およびフェニキア人の数学について丁寧に解説し、読者に対して歴史的な背景と重要な人物について理解を深めていただくことを目指しました。次回は、ギリシャの数学家たちについて詳しく触れることで、数学の発展の流れをさらに追っていきます。

\chapter{Section 2}

\#\#\# 第二章：アテネの数学教師とその時代

\#\#\#\# アナクサゴラス、ソフィスト、そしてヒッパシウス

私たちが学び直す数学の旅を始める前に、古代ギリシャの数学の起源について少し触れましょう。ここから始まるのは、約2500年前にアテネで活躍した数学者たちの時代です。

まず最初に登場するのは、アナクサゴラス（前420年頃）です。彼はソフィストと呼ばれる哲学者や弁士として知られていましたが、数学にも深い関心を持ちました。その中でも特に興味深かったのは3つの問題でした。

1. **特定の半月形の平方化**：これは、半円から作られる特定の形状を正方形に変換するという挑戦的な問題です。
2. **立方体の二倍問題**：これは、既存の立方体の辺を2倍に増やして新しい立方体を作る問題です。この問題は後々、数学の重要な分岐点となることになります。

\#\#\#\# ヒッパシウスとその時代

次に登場するのは、ヒッパシウス（前420年頃）です。彼は特に興味を持っていた3つの問題があります：

1. **文字を幾何図形の描写に使用**：これは、数式や文字を使って幾何学的な図形を表現する方法を確立することです。
2. **帰着法の導入**：これは、複雑な問題を小さな部分に分割し、それぞれを解決することで全体の問題を解くという手法です。
3. **特定の半月形の平方化**：これは、前項で述べたアナクサゴラスが興味を持っていた同じ問題です。

\#\#\#\# プラトンとエウドックスス

プラトン（紀元前429年～348年）は、幾何学における解析法の導入に貢献しました。彼はまた立方体の二倍問題に関する定理を発表しました。一方、エウドックスス（紀元前408年～355年）は黄金分割に関する定理を証明し、さらに「窮竭法」を導入しました。これは、無限小の概念を用いて面積や体積を計算する方法です。

\#\#\#\# メネアクムとその弟子たち

メネアクム（紀元前340年頃）はプラトンとエウドックススの弟子で、放物線、双曲線、楕円線についての議論を行い、立方体の二倍問題に対する2つの解法を提案しました。彼の弟子たちはアリストオイデュスやテアエトetusなど、それぞれ独自の貢献を果たしました。

\#\#\#\# アリストテレスとその時代

最後に、アリストテレス（紀元前384年～322年）は機械学に関する質問で有名ですが、同時に量を示すために使用される文字の導入にも貢献しました。彼の著作は数学だけでなく、哲学や物理学など幅広い分野に及ぶ重要な資料となっています。

\#\#\# 第IV章：最初のアレクサンドリア学派（紀元前300年～前30年頃）

\#\#\#\# 埃及学派の基礎

この章では、アレクサンドリア学派がどのように発展したかを追います。アレクサンドリアは、数学や科学の中心地となり、多くの数学者たちが集まり、新たな知識を創造しました。

1. **オイキドス・ユシップス・ニコテレス**：彼は紀元前3世紀に活躍し、特に平面幾何に関する著作を残しています。
2. **アリストアクス（紀元前310年頃～250年頃）**：彼は太陽からの距離を決定する方法についての研究を行いました。
3. **コノン・ディソティオス**：彼は幾何学と天文学に関する重要な著作を残しています。

\#\#\#\# アーギメネデスとその時代

アーギメネデス（紀元前287年～212年）は、平面幾何と三次元幾何に関する多くの著作を残しました。彼の代表作には「牛の問題」が含まれており、算術と幾何学の深い関連性を示す重要な例です。

\#\#\#\# アポロニウスとその時代

アポロニウス（紀元前260年頃～200年頃）は、交点に関する著作を残し、立方体の倍積問題の解決に貢献しました。彼の著作は、現代でも高く評価されています。

\#\#\#\# エラトステネスとその時代

エラトステネス（紀元前275年～194年）は、「篩」で有名ですが、その他の数学的貢献も重要です。彼は三角法の基礎を確立し、科学的な天文学者としての地位を築きました。

\#\#\# 第五章：第2アレクサンドリア学派（前30年～後641年）

\#\#\#\# チリスト後1世紀

この時期には、セレヌス・メネラウスが活躍しました。彼は特に三角法の発展に貢献し、中世ヨーロッパにおける算術の進展にも影響を与えました。

\#\#\#\# セレヌス・メネラウス

セレヌス（紀元後10年頃～75年頃）は、三角法の発展に重要な役割を果たしました。彼の著作は、後の数学家たちにとって非常に価値のある資料となりました。

\#\#\#\# ニコマコスとその時代

ニコマコス（紀元後100年頃）は、中世ヨーロッパにおける算術の進展に貢献しました。彼の著作は、算術と代数の基本問題を扱っています。

\#\#\#\# チリスト後2世紀

この時期には、スミルナのテオン（168年に没）が活躍しました。彼は『アルマゲスト』という重要な天文書を編纂し、また幾何学にも貢献しました。

\#\#\#\# テオンとその弟子たち

テオンの弟子であるタイマラディスも重要な数学家でした。彼らの著作は、後の数学発展に大きな影響を与えました。

\#\#\# まとめ

この章で紹介した数学者たちは、それぞれ独自の視点から数学を進化させました。彼らの努力が積み重なって今日の数学があることを思い起こすと、感動しますね。次回はさらに詳しい話題に取り組んでいきましょう。

では、この旅を一緒に楽しみながら学び直していきましょう。あなたの新たな数学の冒険が始まります！

\chapter{Section 3}

\#\#\# 第十一世紀：モロッコの教師たちとアルジャベ

\#\#\#\# ジバー・イブン・アフラとアルザチェル

こんにちは、皆さん。今日は第十一世紀の数学を学びましょう。この時代には、モロッコで活躍した数学者たちが重要な貢献をしました。

まず最初に、ジバー・イブン・アフラ（Jabir ibn Aflah）について触れましょう。彼は1050年頃から1125年頃まで活動し、数学と天文学の分野で多くの著作を残しています。特に、彼は三角法の発展に貢献しました。三角法とは、三角形の角度や辺の関係性を研究する数学の一分野です。

ジバー・イブン・アフラは、当時のヨーロッパで広く知られるようになった「アルジャベ」（Al-Jabr）という概念に大きな影響を与えました。アルジャベとは、「補完」という意味を持ち、方程式を解くための手法を指します。

次に、アルザチェル（Al-Zarqali, Al-Azarchi）について触れましょう。彼は1029年頃から1100年頃まで活動し、天文学と数学の分野で活躍しました。アルザチェルは、太陽や月の位置を正確に測定するための装置を開発し、その結果を「アルジャベ」（Al-Zarqali）という名前で知られる著作にまとめました。

これらの数学者たちが取り組んだ三角法とアルジャベは、後の数学家たちにとって重要な基礎となりました。彼らの研究は、ヨーロッパの数学界に大きな影響を与え、アラビア数字や代数学の導入につながる一歩を踏み出しました。

\#\#\#\# アラビア数字の商業への導入

第十一世紀には、アラビア数字が商業活動にも広く採用されるようになりました。この時代に活躍した数学者たちが、アラビア数字の優れた特徴を認識し、その使用を推奨しました。

例えば、ピサのレオナルド（Leonardo of Pisa, 1175年頃～1230年）は、『 Liber Abaci』（算術の書）という著書で、アラビア数字の優れた特徴を詳しく説明しました。彼は、アラビア数字が計算の簡便さと精度を大幅に向上させることを示し、ヨーロッパの商人たちにその利点を伝えたのです。

レオナルドの『 Liber Abaci』は、アラビア数字の導入に大きな役割を果たしました。彼の著作は、数学が日常生活や商業活動にも深く関わるという概念を広め、ヨーロッパの数学界に大きな影響を与えました。

\#\#\#\# 数学の進歩と歴史的背景

第十一世紀から十二世紀にかけて、数学は急速に発展しました。この時代には、多くの数学者たちが三角法や代数学の分野で重要な貢献をしました。彼らの研究は、後の数学家たちにとって重要な基礎となり、ヨーロッパの数学界に大きな影響を与えました。

これらの数学者たちは、当時の社会や科学技術の発展と密接に関連していました。例えば、ジバー・イブン・アフラやアルザチェルは、天文学や地理学の研究にも貢献しました。また、レオナルドのような数学者は、商業活動にも数学を応用し、その成果を広く伝えることで、数学が日常生活にも深く関わるという概念を広めました。

この時代の数学家たちの努力は、ヨーロッパの数学界に大きな影響を与え、後の数学の発展につながりました。彼らの研究は、現代の私たちにとっても貴重な歴史的遺産となっています。

以上で第十一世紀の数学を学びました。次回は十二世紀の数学家たちについて詳しく見ていきましょう。お楽しみに！

\chapter{Section 4}

\#\#\# 第17章 レーニッツと18世紀前半の数学者

\#\#\#\# レーニッツおよびベルヌイ一家

\#\#\#\#\# レーニッツ、1646-1716
ニュートンと同じ時代に活躍した数学の巨匠、ゲオルク・レーニッツは、哲学と文学にも卓越した才能を持ち合わせていました。彼の最大の功績は微分積分学の発展であり、この分野での論争を巻き起こしました。

\#\#\#\#\# 微分積分の起源に関する論争
レーニッツは17世紀末から18世紀初頭にかけて、微分積分学の基本的な概念を開発しました。彼が独自に導いたこれらの概念は、ニュートン自身も独立して開発していたと主張し、この事実をめぐる論争は数学史において大きな話題となりました。

\#\#\#\#\# 無限小解析に関する彼の著作
レーニッツは無限小解析という分野でも多くの貢献を行いました。彼の著作には「Various mechanical problems」があります。この論文では、複雑な機械的な問題を解くための新しい手法を提示しています。

\#\#\#\#\# その作品の特徴
レーニッツの作品は、理論的な深さと実用性に優れていました。彼は常に数学の応用可能性を探求し、それが科学や技術にどのように貢献できるかを考えました。

\#\#\#\# ジェームズ・ベルヌイ、1654-1705
\#\#\#\#\# ヨハネス・ベルヌイ、1667-1748
ジェームズとヨハネスは兄弟で、二人とも数学の巨匠でした。彼らは微分方程式や無限級数に関する重要な発見を行い、これらの成果は現代の数学にも大きな影響を与えています。

\#\#\#\#\# 若きベルヌイ一家
兄弟2人は若くして天才的な才能を発揮し、数学界に大きな衝撃を与えることになりました。彼らの研究は互いに競い合いながらも協力し、18世紀初頭の数学者たちの中で最も活発な交流を生むこととなりました。

\#\#\#\#\# 大陸における解析の発展
ベルヌイ兄弟の活動はヨーロッパ大陸全体で数学の発展に大きく貢献しました。彼らの研究はフランス、ドイツ、イタリアなど多くの国々で広く受け入れられ、数学の新たな潮流を生み出しました。

\#\#\#\# ルポシタル、1661-1704
\#\#\#\#\# バリニョン、1654-1722. モンモールト. ニコル
\#\#\#\#\# パレント. サーニ. デ・グワ. クラメール、1704-1752
\#\#\#\#\# リカティ、1676-1754. ファゴーノ、1682-1766
\#\#\#\#\# クラロア、1713-1765
\#\#\#\#\# ダランベール、1717-1783
これらの数学者たちは、ベルヌイ兄弟の後を継いで解析学の発展に貢献しました。特にダランベールは二階の偏微分方程式の解法において重要な成果を挙げました。

\#\#\#\# ダニエル・ベルヌイ、1700-1782
\#\#\#\#\# 十八世紀の英国数学者
\#\#\#\#\# デービッド・グレゴリー、1661-1708. ハレー、1656-1742
\#\#\#\#\# ディットン、1675-1715
\#\#\#\#\# ブロック・テイラー、1685-1731
テイラーは「テイラーの定理」を発見し、これは現在でも広く使われる数学的な工具です。彼の物理学的研究も非常に重要で、多くの実験と理論的考察を通じて新たな洞察を得ました。

\#\#\#\#\# コッツ、1682-1716
\#\#\#\#\# デモヴレ、1667-1754. 三角法の発展
デモヴレは三角法の発展に大きく貢献し、その成果は現在でも数学の基礎となっています。彼の著作は三角関数の性質や応用について深く考察しています。

\#\#\#\#\# マクラーレン、1698-1746
マクラーレンは幾何学的な発見を行い、特に流動論に関する著作で知られています。彼の研究は流体力学や熱力学など、現代の物理学にも大きな影響を与えています。

\#\#\# 第十八章 ラグランジュ、ラプラスおよび当世の数学者　約1740年−1830年

\#\#\#\# 時代の数学の特性
この時代には解析と力学の発展が著しく進みました。特にオイラー、ラグランジュ、ラプラストンなど、多くの数学者たちが活躍し、現代数学の基礎を築きました。

\#\#\#\# 解析と力学の発展
\#\#\#\#\# オイラー, 1707--1783
オイラーはこの時代の代表的な数学者であり、彼の著作は非常に多岐にわたります。『無限に関する分析入門』（1748年）、『微分計算の原理』（1755年）、『積分計算の原理』（1768-1770年）など、これらの著作は数学教育と研究において重要な役割を果たしました。

\#\#\#\#\# 力学と天文学に関するオイラーの著作
オイラーは力学と天文学にも多くの貢献を行いました。彼の著書『代数入門』（1770年）では、代数の基本的な概念を分かりやすく説明しています。また、力学と天文学に関する研究も豊富で、これらの成果は現代の物理学にも大きな影響を与えています。

\#\#\#\# ラムベルト, 1728--1777
\#\#\#\#\# ブゾート, 1730--1783. テルメリー, 1749--1811. アルボガスト, 1759--1803
これらの数学者たちは、オイラーの後を継いで解析学と力学の発展に貢献しました。特にブゾートは、偏微分方程式の解法において重要な成果を挙げました。

\#\#\#\# ラグランジュ, 1736--1813
\#\#\#\#\# 各種の主題に関するメモリ
\#\#\#\#\# 『解析力学』, 1788
\#\#\#\#\# 『理論』および『関数計算』, 1797, 1804
\#\#\#\#\# 『代数方程式の解法』, 1798.
ラグランジュは、数学と物理学の両方に大きな影響を与えた数学者です。彼の著作には「Various mechanical problems」があります。また、「解析力学」という重要な本を著し、その内容は現代でも学術的な価値が高く評価されています。

\#\#\#\#\# ラグランジュの著作の特性
ラグランジュの著作は理論的で精密でありながらも、実用性にも優れていました。彼は常に数学と物理学の融合を目指し、その研究は多くの分野に広く影響を与えました。

\#\#\#\# ラプラス, 1749--1827
\#\#\#\#\# 天文学と引力に関するメモリ, 1773--1784
ラプラストンは天文学と引力の研究において大きな成果を挙げました。彼の著作は「理論」および「関数計算」という重要な本で、これらの内容は現代でも数学教育に広く用いられています。

この章では、18世紀前半から約1830年頃までの数学者たちがどのように解析学と力学を発展させたかについて詳しく説明しました。彼らの研究は現代の数学や物理学にも大きな影響を与えています。

\chapter{Section 5}

\#\#\# 新分野の創出

\#\#\#\# 数学のこの世紀についての議論の難しさ
現代の数学研究は非常に複雑で、その全体像を把握するのは難しいです。ここでは、主な分野や重要な発見を概観しますが、完全な一覧ではないことに注意してください。

\#\#\#\# 参考文献
- ガウス (1777-1855)
  - 天文学における調査
  - 電気学における調査
  - 《算術的研究》, 1801年
  - 他の発見

- ラグランジュ、ラプラス、ガウスの比較

- ディリクレ (1805-1859)
  - 数論理論の発展

- アイゼンシュテイン (1823-1852)

- ヘンリー・スミス (1826-1883)

- クーマー (1810-1893)
  - 数論に関する他の著者の注記

- 多重周期関数の理論の発展
  - ア贝尔 (1802-1829) の定理
  - ヤコビ (1804-1851)

- リーマン (1826-1866)
  - 楕円とアベイエル関数に関する他の著者の注記

- ウィエリストラス (1815-1897)
  - 楕円とアベイエル関数に関する最近の著者の注記

\#\#\#\# 関数論
- 高次代数の発展
  - オイラー (1759-1833)
  - アルカン (1768-1822) の複素数の幾何学的解釈
  - Sir ハミルトン (1805-1865) のクォータニオンの導入
  - グラスマン (1809-1877) の非可換代数 (1844年)
  - ボール (1815-1864) とデ・モラン (1806-1871)

- ガロワ (1811-1832) の離散置換群の理論
  - ケイリー (1821-1895)
  - スライヴスター (1814-1897)
  - リーディング (1842-1889) の持続群の理論
  - エルミート (1822-1901)

\#\#\#\# 解析幾何学の発展
- 最近の著者による解析幾何学に関する注記

\#\#\#\# 線形幾何学
- 分析に関する最近の著者の名前

\#\#\#\# 合成幾何学の発展
- シュタイン (1796-1863)
  - スタウト (1798-1867)

- 現代の合成幾何学に関する他の著者の注記

\#\#\#\# 非ユークリッド幾何学の発展
- 平行線に関するユークリッドの公理
  - 双曲幾何、球面幾何
  - 合同図形

\#\#\#\# 数学の基礎
- その科目における仮定

\#\#\#\# 運動学
- 力学の理論的発展（グラフィカル）
  - 力学の理論的発展（解析的）

- 力学に関する最近の著者の注記

\#\#\#\# 理論天文学の発展
- ベッセル (1784-1846)
  - レヴィエール (1811-1877)、アダムス (1819-1892)

- 理論天文学に関する他の著者の注記

\#\#\#\# 最近の発展
- 数学物理学の発展

---

\#\#\# 初期の数学

最初に、数学の歴史はイオニアのギリシャ人学校やその以前の時期までさかのぼることは確実ではありません。その後の歴史は三つの時代に分けられます：第一期はギリシャの影響下にある数学の歴史で、これは第Ⅱ章から第Ⅶ章まで讨论されます；第二期は中世とレネサンス時代の数学で、これは第Ⅷ章から第XIII章まで讨论されます；第三期は近代数学で、これは第XIV章から第XIX章まで讨论されます。

\#\#\#\# イジプトとフェニキアの数学

数学の歴史はイオニアの学校から始まりますが、最初にこの科目に関心を示したギリシャ人は、イジプト人やフェニキア人の以前の研究に大きな依存があったことは疑いの余地がありません。これらの民族の数学的達成についての我々の知識は不完全で部分的に推測に基づいていますが、そのようなものでもここでは簡略にまとめられます。

明確な歴史は次の章から始まります。

\#\#\#\# フェニキアの数学

数の科学についてまず言及すると、フェニキア人のこの分野での学識については確実なことは言えません。トリーアとシドンにおける商業取引の規模は、算術の発展を必要とし、その発展には「科学」という名称が適切に使われる可能性があります。

バビロニアでは連続する一列の整数の平方の数値的な値が記された表が見つかり、これは数の性質が研究されていたことを示唆しています。ストラボによれば、フェニキア人は特に数の科学、航海、天文学に注目し、彼らは近隣の親戚であるカルデア人と大きな商取引を行っていました。

また、バッヒクはフェニキア人がバビロニアで使用する重量と測定法を定期的に提供していたと言っています。カルデア人は確かに算術や幾何学にも一定の注意を払っていたと考えられ、これは彼らの天文計算から示されています。無論、その算術における達成度合いがいかなるものであれ、フェニキア人が同程度に精通していたことはほぼ確実であり、またフェニキア人の知識は希薄なものであってもギリシャ人に伝えられた可能性があります。

早期のギリシャ人はフェニキア人から実用的な算術や計算技術を学んだことが多く、おそらく数のいくつかの性質についても学んだかもしれません。ピタゴラスはフェニキア人の出身であり、ヘロドトゥスによれば（ただし疑問視される）、テレスもその血統でした。

算盤や Swan-pan の几乎普遍的な使用が古代の人々に理論的な算術の知識なしでも加減法を行うことを容易にしたことを述べておくことにします。これらの機器は後ほど第 VII 章で説明するので、ここでは十分に言うのは、これらは十進法における数を具体的な方法で表現し、加減法の結果を得るための単なる機械的な過程を可能にする点です。

---

この解説講義を通じて、大人の読者は数学の歴史と発展について深い理解を得ることができます。各時代の重要な人物や発見を丁寧に紹介し、その背景や影響を生き生きと伝えることで、読者の興味を引き出し、学びの喜びを感じていただけるでしょう。

\chapter{Section 6}

\#\#\# 分数と算術：古代文明の数学的探求

分数に如此多くの注意が向けられた理由は、その扱いが初期には困難だったことにあります。エジプト人およびギリシャ人は、分数をいくつかの分子が1である分数の和として簡略化することで問題を解決しました。この規則の唯一の例外は $\frac{2}{3}$ でした。これがギリシャ人の慣行であり、私たちの紀元6世紀まで続きました。一方でローマ人は、分母を一般的に12とし、分数を大体的にはその数の十二分の一として表しました。\index{角度,!六十進法}バビロンiansは天文学において同じ方法を使用し、定数の分母として60を使いました；そして彼らを通じて、現代の角度単位である1度を60等分する分割法が生まれました。つまり、分子と分母の両方の変化を考える必要を避けた形で問題が解決されました。今日では小数を使用することで、定数の分母を維持することがよくあり、これによりローマ人の慣行に逆戻りしています。

分数を考慮した後、アメスは算術の基本的な過程について的一些例に進みます。\index{乗法, 過程}乗法において彼は繰り返し加算に依存していたように見えます。例えばある数 \(a\) を13倍する必要がある場合、彼はまず2倍にして \(2a\) を得、次にその結果を再度2倍にして \(4a\) を得、さらにその結果を再度2倍にして \(8a\) を得、最後には \(a\), \(4a\), および \(8a\) を加算します。おそらく除法も繰り返し減算によって行われたかもしれませんが、彼が結果に至る過程をほとんど説明していないため確実ではありません。これらの例の後、アメスはいくつかの単純な数値方程式の解法に進みます。例えば彼は「山、その七分の一、全体で十九」と述べ、これはある数を見つけ出すことを意味し、その数とその数の七分の一の合計が19になるようにという意味です。そして彼は答えとして \(16 + \frac{1}{2} + \frac{1}{8}\) を与え、これは正しい答えです。

\#\#\# 代数的記号と算術

算術的な部分では、彼は代数的記号の概念についてある程度の理解を持っていました。未知量は常に「heap」という記号で表されます。加法は時として前進する足のペアで表現され、減法は後退する足のペアまたは矢印の飛行で表現されます。等式は、少しごく微小な角度で回転した＜と＞の上にマイナス記号を重ねたような記号で表されます。

\#\#\# 几何学：土地測量から発展

第二に、幾何学について。幾何学は土地測量から発展したとされていますが、数や計算の研究が始まった時期については、文明国家においては不可欠な知識であるため、その具体的な始まりは困難です。しかし、抽象的な幾何学的論理と実用的な土地測量の規則を区別することは比較的容易です。土地測量の方法は古くから存在していたかもしれませんが、古代の普遍的な伝説によれば、幾何学の起源はエジプトに求められるべきだと主張されていました。ギリシャでは生まれず、土地測量の必要性から発展したと考えるのはより確実です。この推測は、「g\~h」（大地）と「metr'ew」（測る）という言葉の由来からも支持されます。

現在まで残っているギリシャ幾何学者の作品を基に判断する限り、彼らは常に抽象的な科学として幾何学を扱いました。絶対的に真実である定理を求め、少なくとも歴史的な時代においては、単位長さが考慮したいくつかの量と互換的でない場合でもその結果が真実に近いものであるという主張をするでしょう。したがって、この名前は彼らの実践を指すものではなく、ギリシャ人がエジプトから学んだ幾何学における使用方法を示唆している可能性があります。これはまた、ギリシャの伝説と一致し、その内容自体が頼りになるように思われます。ヒロドトゥスによれば、尼罗河の定期的な氾濫（河岸線を流し去り、河川沿いの土地の境界を消滅させ、それに伴う税収価値の増減を引き起こす）は正確な測量システムが不可欠に感じさせ、これにより祭司たちによる系統的な研究が行われたとされています。

\#\#\# エジプトの幾何学

フェニキア人やエジプト人の近隣の民族に対して特別な注目が几何学に向けられた理由は見出せない。ユダヤ人は幾何学にあまり注意を払わなかったという小さな証拠として、彼らの聖書に誤りがある。第一王書第七章二十三節と第二 Chronicus 第四章二節では、円の周長が直径の三倍であると述べられている：バビロニア人は$\pi$も3とみなしていた。\footnote{J. Oppert, \emph{Journal Asiatique}, 1872年8月号および1874年10月号を参照。}幾何学の起源|）

次に、ギリシャ人がエジプトから幾何学の知識を得たと仮定すると、エジプトの幾何学の範囲と性質について議論する必要があります。\footnote{Eisenlohr; Cantor, 章Ⅱ; Gow, 著 75, 76; および E. Weyr の著作 \emph{Die Geometrie der alten Aegypter}, ヴィenna, 1884。}

アヘメスの著作以前に幾何学的な結果が知られていたことは、その前世紀に特定の建物の敷地計画において直角を得るための以下の方法が使用されていたという事実から明らかです。\footnote{エジプト神殿の向き付け}

エジプト人は神殿の正確な向きを非常に重視しており、そのため北と南の線、また東と西の線を正確に得る必要がありました。星が昇り降りる点を地平線上で観測し、その間の中央に平面を設けることで北と南の線を得ることができました。東と西の線はこの直角に垂直な線が必要であり、専門的な「ロープ固定者」が使用されました。これらの人は結び目やマークで区切られた長さ比 $3:4:5$ のロープ $ABCD$ を使用しました。長さ $BC$ は北と南の線に配置され、ノック $P$ と $Q$ は節点 $B$ と $C$ に挿入されました。部分 $BA$（常に張っておく）は peg $P$ 回りに回転させ、同様に部分 $CD$ は peg $Q$ 回りに回転させ、端点 $A$ と $D$ が一致するまでしました。その地点は peg $R$ でマークされました。結果として三角形 $PQR$ を形成し、辺 $RP$, $PQ$, $QR$ の長さ比は $3:4:5$ でした。三角形の角 $\angle P$ は直角であり、線 $PR$ が東と西の線を提供しました。\footnote{インダクティブ幾何学|(}

現在でも実用的なエンジニアリングでは、正確な直角を測定するために類似の方法が常に使用されています。この性質は \Eucn@Euc.!a48@\textsc{i}, 48\ref{Euc.i.48} の特定のケースとして導出できます。そして、エジプト人がこの

---

以上のように、分数や算術、そして幾何学について、中学校・高校レベルの知識で理解しやすく説明しました。歴史的背景も含めて、読者に深く共感できる内容を提供しています。

\chapter{Section 7}

\#\#\# 7.1 古代中国人の数学

\#\#\#\# 前提知識
- 方程式、関数のグラフ、三角比、確率の基礎、微分積分の直観的イメージ

\#\#\#\# 講義内容

**古代中国人の数学**

私たちが知る限り、古代中国人は確実に理解していた幾何学的な定理は非常に限られています。その中でも最も有名なのは「3:4:5」または「1:1:\(\sqrt{2}\)」という辺の比を持つ直角三角形において、斜辺上に描かれた正方形の面積が他の二辺上に描かれた正方形の面積の合計と等しいという定理です。これは「ピタゴラスの定理」の一種であり、古代中国人はこの定理を実用的に使用していました。

さらに、半実験的な方法である重ね合わせによる証明が可能な少数の幾何学的定理も彼らが知っていた可能性があります。例えば、正方形や長方形の面積計算などです。これらの知識は十進法で記録されていましたが、その主な特徴は Swan-パンを使用した計算技術と、結果を書く能力にありました。

\#\#\#\# 中国数学の歴史

中国における理論数学の発展については、16世紀に中国へ赴いた jesuit 神父たちによる誤解から生じたものと考えられます。彼らは中国原本の科学と、到着時に見つけた一般的な見解を区別することができませんでした——後者には13世紀以降に中国へ Came きたアラブ人またはインド人神父たちの著作や教えが基盤になっており、その際に球面三角法の知識も導入されました。また、政府機関の一つとして「数学局」と呼ばれるものがあることを発見し、その機能は帝国全体における数学研究の推進と監督だと誤って考えました。実際には、その役割は毎年一度の暦書の作成に限定されており、その中で記された日付や予測が公共や家庭生活において多くの事項を規制していました。

\#\#\#\# 中国数学の特徴

L. A. Sédillot の研究によると、中国人は算術や幾何学における有限なルールを系統的に分類したり拡張したりする試みはありませんでした。また、彼らは観察結果の原因を説明することもありませんでした。これらの特徴から、ギリシャ人が登場する前の数学的成績に関する知識は非常に限られています。

\#\#\#\# まとめ

ギリシャ人は直接または間接的に算盤の使用法、結果を記録するためのシンボル、そして Rhind パピラウスに含まれるか暗示される程度の数学を学んだと考えられます。おそらくこれが彼らが他の民族から借得した知識の総和でしょう。

\#\#\#\# 次章展望

次六章では、ギリシャの影響下における数学の発展を追跡します。第1期はテラスの教えから始まり、紀元前約600年頃からアレクサンドリアがイスラム教徒によって641年頃に占領されるまで続きます。この時期の特徴は幾何学の発展です。

\#\#\#\# 参考文献

これらの学校の歴史については、G. ロリアが『古代ギリシャにおける正確な科学』（モーナーナ、1893年～1900年）、コンタールが第5章から第8章で、G. J. アラマンが『テレスからユークリッドまでのギリシャ幾何学』（ダブリン、1889年）、ゴウが『ギリシャ数学』（ケンブリッジ、1884年）、C. A. ブレッツネッダーが『ユークリッド以前の幾何学と幾何学者』（ライプツィグ、1870年）で、また部分的にハンケルが亡き後に発表した『数学史』（ライプツィグ、1874年）で論じられています。

\#\#\#\# まとめ

イオニア学派とピタゴラス学派の基礎が築かれて以降、考古的研究や推測から歴史の光に出てきます。430年以前まで THESE 学派の哲学者たちの知識を評価するための資料は非常に限られています。当時書かれた異なる数学的論文の大半は失われてしまい、また325年前後頃に Eudemus によって書かれた数学の歴史についての著作も残されていません（Eudemus はアリストテレスの弟子でした）。幸いにも、450年頃に Euclid の『要素』の前半部に対する解説を書いた Proclus は Eudemus の著作に馴染みがあり、その内容を歴史的な参照で自由に利用していました。また、50年前後頃に書かれた Geminus の『数学総論』の一部も残されており、この中では早期のギリシャ幾何学者が使用した証明方法と後の時代におけるそれらとの比較が述べられています。

\#\#\#\# まとめ

イオニア学派の創始者であるテラスは、最初期のギリシャ数学と哲学の学派の創設者です。彼はギリシャの七賢哲の一人であり、紀元前640年頃にミレトゥスで生まれ、紀元前550年頃同地で没しました。テラスに関する生涯についての資料は、伝承によって伝えられた数少ない逸話程度です。

彼の早期の生涯においては、商業と公務に両方に関与していました。そして、二つの逸話から推測すると、彼はビジネスにおける鋭さを発揮していたことがわかります。例えば、ある商取引で相手が詐欺をしようとしたとき、テラスはその詐欺を見抜き、巧妙に対策を講じたと伝えられています。

これらの逸話から、テラスは数学的な洞察力だけでなく、実践的な知恵も兼ね備えていたことがうかがえます。彼の生涯は、数学と哲学の融合を示す良い例と言えるでしょう。

\chapter{Section 8}

\#\#\# Section 8：古代ギリシャ幾何学の偉業

\#\#\#\# (i) 底角が等しい二等辺三角形は存在する（Eucn@Euc.!a05@\textsc{i}, 5）

プロクルスは、これは別の完全に等しい二等辺三角形を取り、それを反転させて最初のものに重ね合わせることで証明されたと暗示しているように見える。これは一種の実験的な証明である。

この命題は、底角が等しい二等辺三角形が存在することを示しています。例えば、三角形ABCとDEFがあり、∠A = ∠D かつ ∠B = ∠E の場合、これらの三角形は等しいと考えることができます。プロクルスは、このような三角形を反転させて重ね合わせることで、その等しさを証明しようとしたのです。

\#\#\#\# (ii) 二直線が交わると、対頂角は等しい（Eucn@Euc.!a15@\textsc{i}, 15）

テラスはこれがあからさまなことだと考えていたかもしれません。プロクルスは、ユークリッドが初めて厳密な証明を与えたと追記している。

二直線が交わると、対頂角は等しいという命題です。例えば、直線ABとCDが交わり、その交点をOとするとき、∠AOC = ∠BOD かつ ∠AOB = ∠COD となります。テラスはおそらくこの事実が自明だと感じていたかもしれませんが、プロクルスはユークリッドが初めて厳密な証明を与えたと指摘しています。

\#\#\#\# (iii) 底辺と底角が与えられた三角形は一意に決定される（\emph{cf.}\ Eucn@Euc.!a26@\textsc{i}, 26）

これは海上的の船の距離を求めるために適用されたと考えられる。底辺は塔であり、底角は観測によって得られたものである。

この命題は、底辺と底角が与えられた三角形は一意に決定されることを示しています。例えば、塔から海を見下ろすとき、塔の高さ（底辺）とその角度（底角）が与えられれば、船までの距離も一意に定まります。これは航海や測量で非常に重要な概念です。

\#\#\#\# (iv) 等角三角形の辺は比例する（Eucn@Euc.!e04@\textsc{vi}, 4, あるいはむしろ Eucn@Euc.!e02@\textsc{vi}, 2）

これはテラスがエジプトでピラミッドの高さを測定するために使用したとされる。プラUTERSの対話の中で、演者はテラスにこう述べる。「あなたの棒をピラミッドの影の端に置き、太陽の光線を使って二つの三角形を作り、それによってピラミッドの高さは棒の長さに対してピラミッドの影の長さに対する棒の影の長さと比例することを証明した。」

この命題は、等角三角形の辺が比例することを示しています。例えば、太陽光線を使って作る二つの三角形（一つはピラミッド、もう一つは棒）において、これらの三角形の対応する辺は比例します。これはエジプト人にとっても知られていなかったとされ、王であるアマシスはこの抽象的な科学の応用に驚いたという。

\#\#\#\# (v) 任意の直径で二等分される円がある（Eucn@Euc.!e04@\textsc{vi}, 4）

これはテラスが発表した可能性もあるが、最も古くから明らかな事実と認識されていたことだろう。

この命題は、円の任意の直径は円を二等分することを示しています。つまり、円の中心を通る直線（直径）は円を均等に分割します。これは非常に基本的な幾何学的原理であり、古代から広く理解されていました。

\#\#\#\# (vi) 円の直径が円周上の任意の点で描く角は直角である（\index{Eucn@Euc.!c31@\textsc{iii}, 31. ref.\ to}エウクレシス \textsc{iii}, 31）

この定理は、テラスが最も驚異的な幾何学的成績として評価されたものであり、円に直角三角形を内接させた後、不朽の神々に牛を捧げたという記述がある。

この命題は、円の直径が円周上の任意の点で描く角が直角であることを示しています。つまり、円の中心を通る直径と円周上の任意の点を結ぶ線分（弦）が作る角度は直角です。これは非常に重要な幾何学的原理であり、円の性質を理解する上で基本的な概念です。

\#\#\#\# まとめ

これらの命題は、古代ギリシャの幾何学者たちが発見し、証明した重要な数学的概念です。特にテラスは、これらの定理を実用的に応用し、例えばピラミッドの高さ測定や航海での距離計算などに活用しました。

\#\#\#\# 逸話と影響

テラスは天文についても書いた。そして彼の同僚たちの中で、哲学者としての彼の名よりも、幾何学者としての彼の名前の方がより有名だったと言われています。ある話では、夜に散歩しているとき、星を見つめすぎて足元が見えない状態で壕に落ちたという。その老婦人は「自分の足元が見えないのに、天象をどうやって知っているのか」と言った。この逸話はしばしば無実用的な哲学者の特徴を示すために引用された。

天文的な詳細は省きますが、彼は1年が約365日であると教えたことがあり、それ以前には毎月30日ずつ計12ヶ月という計算方法が使われていたとは言われています。彼の前輩たちは季節を適切な位置に保つために時々月を追加したと言われており、その場合、1年は平均して360日を超えることを認識していた可能性があります。

テラスは、彼が予言した時期に太陽食が実際に起こったとされています。実際の日付は紀元前585年5月28日か紀元前609年9月30日のいずれかです。この予言とその実現は彼の教えに非常に高い威信を与え、ギリシャの七賢人の一員として名を残すことに貢献しましたが、おそらく彼はエジプトやカルデアの日食記録の一つを利用しただけだった可能性が高いです。これらの記録によると、太陽食は約18年11日の周期で再発生すると予想されていました。

テラスの弟子にはアナキマンドロス, アナンキメネス, マメンコス、そしてマンドラトゥスがいました。最後に挙げられた三人についてはほとんど不知です。アナキマンドロスは紀元前611年に生まれ、紀元前545年に亡くなり、メシューの学校の頭取としてテレスを継承しました。スイダスによると彼は幾何学に関する論文を書いたとされ、伝説では球体の性質に特に注意を払い、無限という概念が空間や時間における哲学的な思想に関連していることを詳しく説明したと言われています。彼は地球と天体の球体を作製しました。

アナキスマネラスはギリシャにゴノモンやスタイルの使用を導入したとされる。これは原則として、水平な地面に垂直に立てられた棒だけから構成されていた。最初は太陽時計として使われ、その場合、三つの同心円の中点に置かれていたので、2時間ごとに影の先端が各円から次の円へ移動した。ゴノモンはパOMPIーニやツスクラムで発見されている。彼はこれらのスタイルを使って自己の正午時の太陽の高さを特定するために、その日の朝と夕方の影を投げた線を描き、それらが形成する角度の二等分平面を見つけたとされる；その後、一年のどの時期に太陽の正午時の高さが最大で最小であるかを観測することで夏至冬至を得た。さらに、夏至と冬至の正午時の太陽の高さの和の半分から地球の傾斜角度（場所の高度を決定する）、そしてそれらの差の半分から黄道の傾斜角度を見つけた。彼がスパルタの緯度を実際に特定したのは確実だが、その他の天文学的な推論を行ったかどうかは疑問である。

これらの命題と逸話は、古代ギリシャの数学と天文学者たちの偉大な業績を浮き彫りにし、彼らの知的好奇心と探求精神を生き生きと伝えてくれます。

\chapter{Section 9}

\#\#\# Section 9: ピタゴラス派と幾何学

\#\#\#\# ピタゴラス派の再興と発見

ピタゴラス派は、政治的な影響を失いながらも、テランティウムに本部を移し、哲学的・数学的な団体として再興しました。約100年以上にわたって栄え続けました。ピタゴラス自身は著書を出版しなかった。その学派の仮定では、知識は共通して保持され、外部から隠されており、新たな発見の栄光は常にピタゴラスの創設者に帰属すべきとされました。そのため、ヒッパシウス（紀元前470年頃）が誓いを破り公的にその発見を自慢したことで溺死したとされています。団体が徐々に分散するにつれて、この習慣は廃止され、彼らの教えや学説の内容を含む論文が書かれ始めました。フィロラウス（紀元前370年頃）によって作成された著作は、プラトンがそのコピーを得たと伝えられています。

\#\#\#\# ピタゴラス派の教育理念

ピタゴラスは主に道徳改革者および哲学者でしたが、その道徳と哲学の体系には数学的な基礎がありました。彼の数学的研究は、哲学の体系を説明することが教えの主要な目的であったため、道徳や哲学へと導くものでした。ピタゴラシアンは、扱う数学的主題を四分に分けました：絶対的な数または算術、応用された数または音楽、静止する量または幾何学、そして動く量または天文学。この「四分」（quadrivium）は長い間、自由教育のために必要なかつ十分な学習課程として考慮されてきました。

\#\#\#\# 極めて重要な幾何学

まず、その幾何学について。ピタゴラスはおそらく、ユークリッドの『元素』初編と二番目の内容に関する平行線、三角形、平行四辺形に関する主な事項を含むもの、そして無理数に関する一部の基本的な命題を含む少数の孤立した定理について知っていたし、教え込んだ。しかし、彼の証明がすべて正確ではなかったという疑念があり、特に定理の逆命题がしばしば証明なしに前提として使われていたということが推測されます。

ピタゴラスのこの主題に関する全般的な教義を再現することはほとんど不可能ですが、プロクルスによる『元素』の注記と他の著者の数少ない散在した言及から、以下の命題が含まれていたことがわかります。これらは大半が面積の幾何学に関するものである。

(i) これらの命題は定義から始まり、おそらく数学的概念と哲学との関連を述べる声明よりも、使用される用語の説明より多かった。点という概念については、「位置を持つ単位」という定義が保存されている。

(ii) 三角形の角の和は二つの直角に等しいことが示された（『元素』初編32）。その証明は保存されており、この証明では『元素』初編13と初編29の一部の命題の結果が引用されている。その証明はユークリッドのものとほぼ同じであり、最後に述べられた二つの命題の証明もピタゴラス自身によるものである可能性が高い。

\#\#\#\# 直角三角形の性質

Pythagoras は、ユークリッドの \textsc{i}, 47 及び \textsc{i}, 48 で与えられる直角三角形の性質を証明したことは確実です。これらの命題の証明がユークリッド自身の発明であることが知られているため、最初の定理に対するピタゴラスの原始的な証明は何だったのかという興味深い疑問が湧いた。

この定理の証明はおそらく以下の二つの中の一つであったと推測されます（1896年から1899年にかけて発行された \emph{American Mathematical Monthly Journal} の第3巻から第6巻に収録された、Euc.~\textsc{i}, 47 の証明の集約が100種類以上含まれている）。

(\alpha) 辺 $ABCD$ の正方形は、ユークリッドの \textsc{ii}, 4 参照のように、二つの正方形 $BK$ と $DK$ および二つの等しい長方形 $AK$ と $CK$ に分割される：つまり、これは正方形 $FK$ 上の平方と正方形 $EK$ 上の平方と四倍の三角形 $AEF$ の和に等しい。もし点 $G$, $H$, $E$ をそれぞれ辺 $BC$, $CD$, $DA$ 上にとり、$BG$, $CH$, および $DE$ がそれぞれ $AF$ に等しいとすると、容易に示されるように四角形 $EFGH$ は正方形であり、三角形 $AEF$, $BFG$, $CGH$, および $DHE$ は等しい：したがって、正方形 $ABCD$ も正方形 $EF$ 上の平方と四倍の三角形 $AEF$ の和に等しい。従って、正方形 $EF$ 上の平方は正方形 $FK$ 上の平方と正方形 $EK$ 上の平方の和に等しい。

(\beta) 直角三角形 $ABC$ において、$A$ が直角であるとする。$AD$ を $BC$ に対して垂直に引く。三角形 $ABC$ と $DBA$ は相似であるから、

\[
\therefore BC : AB = AB : BD.
\]

同様に

\[
BC : AC = AC : DC.
\]

従って

\[
AB^2 + AC^2 = BC(BD + DC) = BC^2.
\]

この証明は、ユークリッドの \textsc{ii}, 2, と密接に関連しています。

\#\#\#\# 結論

ピタゴラス派は、数学を哲学や道徳に結びつけることで大きな影響を与えました。彼らの幾何学的発見は、現代でも依然として重要な位置を占めています。この章を通じて、私たちが理解したのは、数学者たちがどのようにして知識を追求し、その成果を共有してきたかです。ピタゴラス派の努力と献身は、私たちの数学教育に深く影響を与え続けています。

以上で、ピタゴラス派の幾何学について詳しく説明しました。この内容を通じて、読者の皆さんは再び数学の魅力を感じ取ることができるでしょう。

\chapter{Section 10}

\#\#\# 10章：ピタゴラス派と数学的探求

\#\#\#\# 特別な数と直角三角形

ピタゴラス派は、特別な数 $(2n^2+2n+1)$, $(2n^2+2n)$, および $(2n+1)$ を直角三角形の辺に見事に適用しました。カントールは、ピタゴラスがこの事実を幾何学的命題 Euc. \textsc{i}, 47（現在では有名なピタゴラス定理）を発見する前に知っていたと考えています。

これらの数のより一般的な表現は $(m^2+n^2)$, $2mn$, および $(m^2-n^2)$、またはそれらの倍数です。ピタゴラスが得た結果は、$m=n+1$ を仮定することでこれらの式から導き出せます。その後の時代には、アキチタスとプラトンは $n=1$ として等価な規則を与えました；ディオファントゥスは一般の表現を知っていました。

\#\#\#\# 三角数

ピタゴラス派はさらに、三角数についても深い理解を持っていたようです。三角数は平面に並べたコインの数の合計を表します。下部の行には $n$ が含まれており、各続く行は1つ少ない：したがって、それは系列
\[
n + (n-1) + (n-2) + \ldots + 2 + 1,
\]
すなわち、$\frac{1}{2} n(n+1)$と等しいです。4に対応する三角数は10となります。

有名なルシアンの文章によると、ピタゴラスが商人に教えたのは「数を数える方法」でした。しかし、商人はすでにその方法を知っていたと言います：「それ Already know.」ピタゴラスは、「どのように数えていますか？」と尋ねます。商人は、「一つ、二つ、三つ、四つ—」と答えると、ピタゴラスは止めて言います：「あなたが四と感じるものは十であり、完璧な三角形で私たちのシンボルです。」

ピタゴラス派は、数をその整数の部分因子の合計と比較して分類しました。数を過剰数、完全数、欠陥数と呼び、これらの部分因子の合計が数より大きい場合、等しい場合、または小さい場合に応じて：最初は偶数のみに制限されました。

第三のグループの問題は、比を形成する数についてであり、おそらくこれはユークリッドの第五巻のように幾何学的な助けを使用して議論されたものでした。最後に、ピタゴラス派は算術、幾何学、調和、音楽的進行数列の数列に関心を持ちました。

\#\#\#\# 四部学科と哲学理論

ピタゴラス派が四部学科の応用科目やそれらに基づく哲学理論について知っているのは非常に少ないです。ピタゴラスは、自然界で見られる特定の数値関係に強く影響を受けたように思われます。結晶学の一部に関する彼の知識があったという説が唱えられています。

彼は弦の振動によって発生する音の高さは弦の長さに関連していることを認識していたと考えられます。特に、音、その五度と八度を生成する長さ比は $2:3:4$ で、音楽的な進行の項を形成していました。また、彼は天文学的惑星が地球からどの距離にあるかについても音楽的な進行に従っていると考えており、宇宙体が空間を通過する際に調和した音を発していることにも着目していたと思われます：このため「球の調和」という表現が生まれました。

これらの結論や類似の結論は彼に、宇宙の秩序と調和の説明は数の科学の中に求められるべきであり、数は形態の原因であるだけでなく正確な測定にも欠かせない要素であるという考えを示唆しました。したがって、彼は特定の性質を持つ特定の数と幾何学的図形に特別な性質を与えることに進みました。

\#\#\#\# アーキタス

アーキタス（Archytas）は紀元前400年頃のテラントゥムの有力な市民の一人で、彼は城の長を7回も務めました。その当時の同僚たちに対する影響力は非常に大きく、一度はディオニシオスによってプラトノの命を救うことに使われました。

アーキタスは奴隷や都市内の子供たちの快適さと教育に注意を払ったことで知られていました。テラントゥム近海で船難に遭い、体が岸に流されましたが、より厳格なピタゴラス派の眼には、彼が創設者によって定められた学問の道から逸脱したため、適切な罰と見做されました。

アーキタスは幾何学の助けで機械工学を扱いました。彼は引致索の理論を発明し、研究し、飛行する鳥や他のいくつかの巧妙な機械玩具も考案したとされる。曲線を構築したり問題を解決するために様々な機械装置を導入しました。

プラトノはこれらが幾何学の知的練習としての価値を損なうと反対し、後期のギリシャ幾何学者たちは、直尺と Compass の2種類の道具のみを使用することに制限された。アーキタスはまた天文学にも興味を持ち、地球は24時間で自軸を中心に回転する球体であり、その周りを天体が動いていると教えた。

アーキタスは立方体を二倍にする問題，すなわち与えられた立方体の体積が2倍になる立方体の辺を見つけ出す方法の最初の解決策の一つを与えました。これは古代において最も有名な問題の一つでした。

アーキタスによって与えられた構成は以下のようになります。円柱の底面の直径 $OA$ 上に半円を描き、その平面が円柱の底面に垂直であるようにします。この半円を含む平面が $O$ を通る生成線回転すると、半円によって描かれる曲面は円柱内で複雑な曲線で交わります。この曲線は軸が $OA$ で、半頂角が（例えば）60°の直円錐と交わり、その交点 $P$ において、$OP$ の底面への射影が円柱の半径に対応します。

以上の内容を通じて、ピタゴラス派は数学を通じて自然の秩序を探求し、数と形の関係性を深く理解しようとしました。彼らの探求は、現代でも引き続き私たちに感動を与え、新たな視点を与えてくれます。

\chapter{Section 11}

\#\#\# 11章：古代ギリシャ数学の黎明期

\#\#\#\# オイノピデスとエレア派

学校の歴史は、紀元前500年頃に生まれ、紀元前430年頃に没したオイノピデス\index{Oenopides of Chios}から始まります。彼は主に天文学者として活動していましたが、エジプトで幾何学を学び、2つの重要な問題の解法が伝えられています：与えられた直線上から与えられた外点を通る垂線を引く問題（Euc. \textsc{i}, 12）と、与えられた点において与えられた角と同じ角度を構成する問題（Euc. \textsc{i}, 23）。これらの問題は、当時の数学の重要な発展に貢献しました。

イタリアのエレアにも重要な中心地がありました。これはシチリアでゼノファネスによって設立されました。彼の後にはパルメニデス\index{Parmenides}、ゼノン\index{Zeno}、メリッソス\index{Melissus}が続きました。\emph{エレア派}のメンバーは無限級数を用いる必要がある問いに対する困難性を提起し、有名な例として知られるアキレスと亀のパラドックスが存在します。このパラドックスは、彼らの中でも最も目立つメンバーであるゼノンによって提唱されました。ゼノンは紀元前495年頃に生まれ、紀元前435年頃にエレアで反乱の罪で処刑されました。彼はパルメニデスの弟子であり、紀元前455年頃から紀元前450年頃までアテネを訪れたことがあります。

ゼノンは、アキレスが亀の10倍速く走っても、仮に亀が1000ヤードの差をつけてスタートした場合、追いつけることはないと主張しました。なぜなら、アキレスが1000ヤード走ったときには、亀は彼の前進に100ヤードほど残っているはずです。そして、アキレスがその100ヤードを走破すると、またしても10ヤードほど先にいます。このようにして無限に続くと、アキレスは亀に近づく一方で、追いつくことはないでしょう。この論理の誤りは通常、追いつくためには時間も無限に分割されると主張し、その部分は幾何級数的に小さくなるが、それらの合計は有限な時間が得られることで説明されます。その時間後にはアキレスが亀の前にいるはずです。

おそらくゼノンは、この議論は空間が無限に分割可能であるという前提に基づいていると反論し、これは争点となっている問題です。彼自身は量が無限に分割可能ではないと主張していました。

これらのパラドックスはギリシャ人を無限小の使用について疑念を持たせ、最終的には窮竭法の発明につながりました。

\#\#\#\# アテネの数学研究

アテネは紀元前5世紀末頃から数学的研究の中心地となりました。そのような状況が形成された理由はいくつかありました。その世紀には、商業と同盟国からの貢献品を利用することで、アテネはギリシャで最も富饒な都市となりました。また、政治家たちの才能により、半島全体の政治中心地としてアテネが位置づけられました。さらに、アテネの政治優越権を主張するどの国家も、彼女に対する知的優越性は認められていました。その世紀には、ある時期にアテネで主要な思考家たちの一人または複数が代表して参加した思想学校が存在していましたし、アジア小アジアとグレギア大公国で熱心に研究されている新しい科学のアイデアも、様々な機会にアテネの人々に提示されました。

\#\#\#\# アナックスガラスとソフィスト

アナックスガラス\index{Anaxagoras of Clazomenae}は、アテネに居住し、アテネ学派の前駆者となった哲学者たちの中で最も重要な人物として挙げられます。彼はイオニア学派の最後の哲学者の一人でした。彼は紀元前500年頃に生まれ、紀元前428年に亡くなりました。紀元前440年前後アテネに定住し、イオニア哲学の結果を教えることになりました。彼もまたその学派の他のメンバーと同様に天文に大きな関心を持ちました。\index{Astronomy.!z12@Early Greek theories of}彼は太陽がペロポネソスよりも大きいという主張を行い、いくつかの物理現象を神の直接的な作用によるものではなく説明しようと試みたことから、不敬罪で訴追され、有罪判決を受けました。牢獄にいたときには円周率の正方形に関する論文を書いたとされています。

ソフィストはアテネ学派に属するものとは考えられず、アナックスガラスも同様です。しかし彼らはアテネ学派が出現する直前に存在し、準備作業を行ったため、彼らについて数行を費やす価値があります。アテネでの公職の成功には良い演説能力が必要であり、都市の財政力と権力が増大したことに伴い、多くの「ソフィスト」がその地に定住し、議論や修辞学などの教養を教えました。幾何学も生徒の教育の一環として行われることがありました。通常二人のソフィストは特に幾何学を研究しました---これはエリスのヒッパシウスとアントピフロスでした。一方、天文を専門とした人物がいました---メトノン\index{Meton}です。\index{Eleatic School|)}

\#\#\#\# ヒッピアスとその発明

ヒッピアスは最初の幾何学者で、エリス出身のヒッピアス（紀元前420年頃）は算術の専門家として記述されており、しかし彼はその曲線「四分線」という名前の発明により最も知られています。この曲線を使用して角度を3等分したり、あるいは任意の比率で分割することができます。「角度の3等分」、「四分線」、「円の平方」。円の半径が中心Oから位置OAを通じて直角まで回転し、同じ時

---

この解説講義では、古代ギリシャ数学の重要な人物や概念について詳しく説明しました。オイノピデスとゼノンのパラドックスは、無限小の使用に関する疑問を提起し、最終的には窮竭法の発明につながりました。また、アテネの数学研究の中心地としての地位や、アナックスガラス、ソフィスト、ヒッピアスなどの人物についても詳しく触れました。これらの知識は、現代の数学の基礎を理解する上で重要な役割を果たします。

\chapter{Section 12}

\#\#\# ユドックススとその時代の数学家たち

ユドックススはアテネの数学者たちの中で最も優れた者であり、キザキウスの学校を創設したとされています。この学校とアテネの学校との間には非常に密接な関係があり、両者の歴史を区別することが難しいほどです。ヒポクラテス、プラトン、テアテトゥスはアテネの学校に属していたという一方で、ユドックスス、メナエクムス、アリストアイウスはキザキウスの学校に属していたとされています。

両者の最初期メンバーはアーキトゥスまたはその弟子テオドロス・カイレヌスの影響を受け、主題に対する取り扱いにも違いがなく、一緒に扱うことが適切です。これらの学校の幾何学者たちは、特に興味を持っていた3つの問題について取り組んでいました：

1. **立方体の倍積**：与えられた立方体の体積を2倍にするための新しい立方体の辺の長さを決定する。
2. **角の三分割**：特定の角度を3つに等分する方法を見つける。
3. **円の平方**：与えられた円の面積と等しい面積を持つ正方形の辺の長さを決定する。

これらの問題は、三次方程式の解を求めることに相当します。しかし、直線と円のみを使用した構成では無解であることが判明しました。双曲線を用いることで、これらの問題は解決できるようになりました。第三の問題については、リンデマンによってその比が有理代数方程式の根でないことを示しており、ユークリッド幾何学では無解です。

アテネ人とキザイア人はこれら3つの問題を解決することができませんでしたが、試みは多くの新しい定理や手続きの発見につながりました。その後期プラトン派は、当時の全ての幾何学的定理を集積し、系統的に整理しました。これらの集積は、ユークリッドの『要素』（書籍 \textsc{i--ix}, \textsc{xi}, および \textsc{xii}）に含まれる命題の大半を含み、また双曲線に関するより基本的な定理もいくつか含まれていました。

\#\#\# ヒポクラテスとその貢献

ヒポクラテスは、ギリシャの最高峰の幾何学者の一人でした。紀元前470年頃キシオスで生まれ、商売人として生活を送りました。彼がアテネの裁判所に訴訟を起こした際には、自分の財産を取り戻すために自己の幾何学学校を開いたとされています。

ヒポクラテスは最初の初等幾何学教科書を著し、おそらくユークリッドの『要素』はその基礎に築かれたものであると言えます。彼は「点Aが立っている点」や「線ABが示されている点」といった表現を使い、図形内に文字を使用することを初めて導入したとされています。

ヒポクラテスは線の平方を「\mytextgreek{d'unamis}」という言葉で表し、この言葉に代数における「冪」の技術的な意味を与えた。彼は直角三角形の各辺の平方の総和が斜辺の平方と等しいという命題を、「直角三角形の各辺の冪の合計は斜辺の冪と等しい」という形式で述べたと言われています。

\#\#\# ヒポクラテスの証明手法

この教科書において、ヒポクラテスは一つの定理を別の定理に「還元」する方法を導入しました。その方法が証明されれば、提案されたことが必然的に成り立つというものです。この手法は反証法（矛盾還元法）を示すものであり、おそらく以前にも時々使用されていましたが、ヒポクラテスはそれを合法的な証明の形態として指摘し、その応用可能性を広く示しました。

円の幾何学について彼は詳細に論じ、相似な円弧は等しい角度を含むことを証明しました。また、円の中接弦によって作られる角は、その円弧が半円より小さいか、等しいか、大きいかにより、直角よりも小さく、等しく、大きくなることも証明しました（Euc. Ⅲ, 31参照）。おそらく彼はユークリッドの第三巻に含まれる他のいくつかの命題も確立しました。相似な円は直径の平方比に比例すること（Euc. Ⅻ, 2参照）と、相似な円弧は弦の平方比に比例することも確立したと思われます。

\#\#\# ヒポクラテスの正方形化問題

ヒポクラテスの最も有名な発見は、円の正方形化と立方体倍積に関連していました。彼の影響により、これらの問題はアテネ学派の歴史上重要な役割を果たしました。

次に示す命題は、彼が正方形化問題に取り組んだ方法を充分に説明します。

\begin{center}
\labelincludegraphics{061a}
\end{center}

彼はまず、半円と同じ弦を底とする四分之一弧で囲まれた月面の面積を求めることから始めました。これは次のようにしました。直角二等辺三角形 \(ABC\) を半円 \(ABOC\) に内接させており、その中心は \(O\) です。\(AB\) と \(AC\) を直径として半円を描きます（図 \ref{061a} 参照）。すると、ユークリッドの第1巻、47命題により
\[

\end{document}
